\section*{Capítulo \ref{ch:b2:amines} --- Aminas y compuestos nitrogenados}

\begin{solution}{exo:b2:amines:1}
Primaria; secundaria; terciaria; sal de amonio cuaternario; primaria (\omterm{def:b2:huckel:aromatic}{aromática}).
\end{solution}

\begin{solution}{exo:b2:amines:2}
Etanamida $<$ anilina $<$ amoníaco $<$ metilamina ($\mathrm pK_a$ de los ácidos conjugados
$-0.6$, 4.60, 9.26, 10.66).
\end{solution}

\begin{solution}{exo:b2:amines:3}
\textit{N}-Metilciclohexanimina, \ce{C6H10=N-CH3}: amina y cetona en medio débilmente ácido
(pH 4--5), eliminando el agua formada (agente desecante o destilación azeotrópica).
\end{solution}

\begin{solution}{exo:b2:amines:4}
Bromobenceno; benzonitrilo; fenol; 4-hidroxiazobenceno \ce{C6H5-N=N-C6H4-OH} (acoplamiento en para respecto al
grupo \ce{O-} del fenóxido).
\end{solution}

\begin{solution}{exo:b2:amines:5}
$\mathrm pK_a = 9.80$: a pH 7 la razón $[\ce{(CH3)3NH+}]/[\ce{(CH3)3N}] = 10^{2.8} \approx 630$: predomina la
forma protonada; el zumo de limón (pH cercano a 2) la protona prácticamente toda. Agitar una
disolución orgánica de una amina con ácido acuoso diluido transfiere la amina, en forma de ion amonio,
al agua; una base la libera de nuevo.
\end{solution}

\begin{solution}{exo:b2:amines:6}
1-(Ciclohex-1-en-1-il)pirrolidina. Tras la adición y la pérdida de agua, el nitrógeno lleva una
carga positiva (un ion iminio) y ningún hidrógeno que perder; se pierde en cambio un protón del carbono
contiguo, que forma el enlace \ce{C=C}.
\end{solution}

\begin{solution}{exo:b2:amines:7}
Benzaldehído con propan-2-amina, o propanona con bencilamina, cada pareja con \ce{NaBH3CN} a pH
cercano a 6.
\end{solution}

\begin{solution}{exo:b2:amines:8}
Propiofenona (1-fenilpropan-1-ona), \ce{C6H5COCH2CH3}: una sola adición da el anión \omterm{def:b2:amines:imine}{imina},
que se hidroliza a la cetona.
\end{solution}

\begin{solution}{exo:b2:amines:9}
Ftalimida potásica + 1-bromohexano: \textit{N}-hexilftalimida; después la hidrazina (o una hidrólisis)
libera la hexan-1-amina. El nitrógeno de la ftalimida alquilada no tiene par no enlazante disponible (está
deslocalizado sobre dos carbonilos) ni le queda ningún \ce{N-H}: no puede alquilarse una segunda vez.
\end{solution}

\begin{solution}{exo:b2:amines:10}
El agua de bromo broma la anilina en las tres posiciones orto y para respecto al grupo
amino: 2,4,6-tribromoanilina. La \omterm{def:b2:amines:diazonium}{diazotación}, seguida de \ce{H3PO2}, sustituye el grupo amino por H: los tres
bromos quedan entonces en posición 1,3,5 entre sí.
\end{solution}

\begin{solution}{exo:b2:amines:11}
4-Nitroanilina + \ce{NaNO2} + \ce{HCl} a 0--\qty{5}{\degreeCelsius}: cloruro de
4-nitrobencenodiazonio. Acoplamiento con fenol en disolución débilmente básica (fenóxido, fuertemente activado) en la
posición para respecto a \ce{O-}: 4-[(4-nitrofenil)diazenil]fenol, \ce{O2N-C6H4-N=N-C6H4-OH}, un colorante
naranja.
\end{solution}

\begin{solution}{exo:b2:amines:12}
El grupo trimetilamonio es un grupo saliente malo y voluminoso, y hace más ácidos y más accesibles los hidrógenos en $\beta$ del
lado del \ce{CH3}; la base arranca un hidrógeno del carbono menos
sustituido: but-1-eno, el alqueno menos sustituido (producto de Hofmann).
\end{solution}

\begin{solution}{pb:b2:amines:1}
\textbf{1.} $\mathrm{{}^-O_3S{-}C_6H_4{-}NH_3^+}$: el grupo sulfónico, ácido, ha protonado la amina.
\textbf{2.} El ion dipolar es poco soluble en agua; el carbonato desprotona el grupo amonio,
dando el sulfanilato de sodio, soluble.
\textbf{3.} \ce{NaNO2 + HCl -> HNO2 + NaCl}.
\textbf{4.}
\[
  \mathrm{{}^-O_3S{-}C_6H_4{-}NH_2 + HNO_2 + H^+ \to {}^-O_3S{-}C_6H_4{-}N_2^+ + 2\,H_2O} .
\]
\textbf{5.} La sal de diazonio se descompone (en el fenol y \ce{N2}) cuando se calienta; el ácido nitroso
también se descompone.
\textbf{6.} $5.00/173.19 = \qty{28.9}{mmol}$: \qty{28.9}{mmol} de \ce{NaNO2}, \qty{1.99}{g}.
\textbf{7.} El ácido nitroso en exceso reaccionaría con el compañero de acoplamiento; se detecta con papel de almidón--yoduro
y se destruye con urea o con ácido sulfámico.
\textbf{8.} Su anillo está fuertemente activado por el grupo dimetilamino, un dador potente.
\textbf{9.} En para respecto al \ce{N(CH3)2}: orientación orto/para, y las posiciones orto están impedidas por los
grupos metilo.
\textbf{10.} $\mathrm{{}^-O_3S{-}C_6H_4{-}N{=}N{-}C_6H_4{-}N(CH_3)_2}$.
\textbf{11.} En ácido fuerte el grupo dimetilamino se protona: \ce{-NH(CH3)2+} desactiva el anillo
y el acoplamiento se detiene.
\textbf{12.} El producto se forma en su forma ácida protonada, en parte soluble; en medio básico el sulfonato
de sodio, menos soluble en la disolución salina, precipita.
\textbf{13.} Su carga positiva está deslocalizada sobre el anillo y el nitrógeno terminal es un centro
electrófilo pobre; solo reaccionan los anillos activados por \ce{-O-}, \ce{-OH} o \ce{-NR2}.
\textbf{14.} Amarillo anaranjado en medio básico, rojo en medio ácido.
\textbf{15.} Sobre un nitrógeno del grupo azo (la forma protonada se estabiliza por conjugación con
el grupo dimetilamino).
\textbf{16.} Hacia $\mathrm pK_a \pm 1$: aproximadamente de 3.1 a 4.4, en torno a $\mathrm pK_a = 3.47$.
\textbf{17.} Dos anillos y el grupo azo forman un sistema conjugado largo: la separación entre el nivel $\pi$ ocupado
más alto y el vacío más bajo es pequeña (\cref{prop:b2:huckel:gap}) y la absorción
cae en el visible.
\textbf{18.} La protonación cambia la distribución de carga en el sistema $\pi$ y estrecha aún más la
separación: la absorción se desplaza a longitudes de onda mayores, y el color pasa del amarillo al rojo.
\textbf{19.} \qty{173.19}{g/mol}.
\textbf{20.} \qty{327.33}{g/mol}.
\textbf{21.} \qty{28.87}{mmol}.
\textbf{22.} $0.02887 \times 327.33 = \qty{9.45}{g}$.
\textbf{23.} La descomposición de la sal de diazonio, un acoplamiento incompleto, la solubilidad del producto en
las aguas madres y en los lavados.
\textbf{24.} $0.80 \times 9.45 = \boldsymbol{\qty{7.56}{g}}$ de naranja de metilo.
\end{solution}
